\documentclass[../../main.tex]{subfiles}
\begin{document}
\section{嵌入}\label{section:ch04-嵌入}

有一类浸入特别重要。设 $M$ 与 $N$ 为带边或不带边的光滑流形。$M$ 到 $N$ 的一个\textbf{光滑嵌入}是一个光滑浸入 $F:M\to N$,它同时还是拓扑嵌入,即在子空间拓扑下是到其像 $F(M)\subseteq N$ 上的同胚。光滑嵌入既是拓扑嵌入又是光滑浸入,而不仅仅是恰好光滑的拓扑嵌入。

\begin{exercise}\label{exercise:ch04-embedding-composition}
证明光滑嵌入的复合仍是光滑嵌入。
\end{exercise}

\begin{example}[Smooth Embeddings]\label{example:ch04-smooth-embeddings}
若 $M$ 是带边或不带边的光滑流形, $U\subseteq M$ 是开子流形,则包含映射 $U\hookrightarrow M$ 是光滑嵌入。

若 $M_1,\dots,M_k$ 为光滑流形, $p_i\in M_i$ 为任意取定的点,则每个由
\begin{align*}
\iota_j(q)=(p_1,\dots,p_{j-1},q,p_{j+1},\dots,p_k)
\end{align*}
给出的映射 $\iota_j:M_j\to M_1\times\cdots\times M_k$ 都是光滑嵌入。特别地,把 $\bigl(x^1,\dots,x^n\bigr)$ 映为 $\bigl(x^1,\dots,x^n,0,\dots,0\bigr)$ 的包含映射 $\mathbb{R}^n\hookrightarrow\mathbb{R}^{n+k}$ 是光滑嵌入。

习题 4-12 表明: \cref{example:ch04-submersions-immersions}(d) 中的映射 $X:\mathbb{R}^2\to\mathbb{R}^3$ 可下降为环面 $S^1\times S^1$ 到 $\mathbb{R}^3$ 的光滑嵌入。
\end{example}

为了更充分地理解一个映射是光滑嵌入意味着什么,记住一些单射光滑映射却不是光滑嵌入的例子是有益的。下面三个例子说明这种情况发生的三种颇为不同的方式。

\begin{example}[A Smooth Topological Embedding]\label{example:ch04-smooth-topological-embedding}
由 $\gamma(t)=(t^3,0)$ 给出的映射 $\gamma:\mathbb{R}\to\mathbb{R}^2$ 是光滑映射且是拓扑嵌入,但它不是光滑嵌入,因为 $\gamma'(0)=0$。
\end{example}

\begin{example}[The Figure-Eight Curve]\label{example:ch04-figure-eight}
考察由
\begin{align*}
\beta(t)=(\sin 2t,\sin t)
\end{align*}
定义的曲线 $\beta:(-\pi,\pi)\to\mathbb{R}^2$。它的像是平面上一个看起来像八字形的集合(\cref{figure:ch04-figure-eight}),有时称为双纽线。(它是满足 $x^2=4y^2(1-y^2)$ 的点 $(x,y)$ 之轨迹,读者可自行验证。)容易看出 $\beta$ 是单射光滑浸入,因为 $\beta'(t)$ 永不消失;但它不是拓扑嵌入,因为其像在子空间拓扑下是紧的,而其定义域不是。
\end{example}

\begin{figure}[ht]
\centering
\includegraphics[width=0.8\textwidth]{fig-04-3.png}
\caption{例 4.19 的8字形曲线}\label{figure:ch04-figure-eight}
\end{figure}

\begin{example}[A Dense Curve on the Torus]\label{example:ch04-dense-curve}
令 $\mathbb{T}^2=S^1\times S^1\subseteq\mathbb{C}^2$ 表示环面,并令 $\alpha$ 为任一无理数。由
\begin{align*}
\gamma(t)=\bigl(e^{2\pi\mathrm{i}t},e^{2\pi\mathrm{i}\alpha t}\bigr)
\end{align*}
给出的映射 $\gamma:\mathbb{R}\to\mathbb{T}^2$ 是光滑浸入,因为 $\gamma'(t)$ 永不消失。它也是单射,因为 $\gamma(t_1)=\gamma(t_2)$ 蕴涵 $t_1-t_2$ 与 $\alpha t_1-\alpha t_2$ 都是整数,而除非 $t_1=t_2$,这是不可能的。

考察集合 $\gamma(\mathbb{Z})=\{\gamma(n):n\in\mathbb{Z}\}$。由 Dirichlet 逼近定理(见下)可知:对每个 $\varepsilon>0$,存在整数 $n,m$ 使 $|\alpha n-m|<\varepsilon$。利用对 $t_1,t_2\in\mathbb{R}$ 有
\begin{align*}
|e^{\mathrm{i}t_1}-e^{\mathrm{i}t_2}|\leqslant |t_1-t_2|
\end{align*}
(因为从 $e^{\mathrm{i}t_1}$ 到 $e^{\mathrm{i}t_2}$ 的直线段短于长度为 $|t_1-t_2|$ 的圆弧),我们有
\begin{align*}
|e^{2\pi\mathrm{i}\alpha n}-1|=|e^{2\pi\mathrm{i}\alpha n}-e^{2\pi\mathrm{i}m}|\leqslant |2\pi(\alpha n-m)|<2\pi\varepsilon.
\end{align*}
因此
\begin{align*}
|\gamma(n)-\gamma(0)|=\bigl|(e^{2\pi\mathrm{i}n},e^{2\pi\mathrm{i}\alpha n})-(1,1)\bigr|=\bigl|(1,e^{2\pi\mathrm{i}\alpha n})-(1,1)\bigr|<2\pi\varepsilon.
\end{align*}
于是 $\gamma(0)$ 是 $\gamma(\mathbb{Z})$ 的极限点。但这意味着 $\gamma$ 不是到其像上的同胚,因为 $\mathbb{Z}$ 在 $\mathbb{R}$ 中没有极限点。事实上,不难证明像集 $\gamma(\mathbb{R})$ 在 $\mathbb{T}^2$ 中实际上是稠密的(见习题 4-4)。
\end{example}

上例与习题 4-4 依赖于下面这个数论中的初等结果。

\begin{lemma}[Dirichlet's Approximation Theorem]\label{lemma:ch04-dirichlet}
给定 $\alpha\in\mathbb{R}$ 与任一正整数 $N$,存在整数 $n,m$ 满足 $1\leqslant n\leqslant N$,使 $|n\alpha-m|<1/N$。
\end{lemma}

\begin{proof}
对任一实数 $x$,令 $\{x\}=x-\lfloor x\rfloor$,其中 $\lfloor x\rfloor$ 是不超过 $x$ 的最大整数。由于 $N+1$ 个数 $\{\{i\alpha\}:i=0,\dots,N\}$ 全都落在区间 $[0,1)$ 中,由抽屉原理,必存在整数 $i,j$ 满足 $0\leqslant i<j\leqslant N$,使 $\{i\alpha\}$ 与 $\{j\alpha\}$ 同落在 $N$ 个子区间 $[0,1/N),[1/N,2/N),\dots,[(N-1)/N,1)$ 之一中。这意味着 $|\{j\alpha\}-\{i\alpha\}|<1/N$,故可取 $n=j-i$ 且 $m=\lfloor j\alpha\rfloor-\lfloor i\alpha\rfloor$。
\end{proof}

下面的命题给出单射浸入成为嵌入的几个简单充分判别法。

\begin{proposition}\label{proposition:ch04-injective-immersion-embed}
设 $M$ 与 $N$ 为带边或不带边的光滑流形, $F:M\to N$ 是单射光滑浸入。若下列任一条件成立,则 $F$ 是光滑嵌入。
\begin{enumerate}[(a)]
\item $F$ 是开映射或闭映射。
\item $F$ 是固有映射。
\item $M$ 紧。
\item $M$ 边界为空且 $\dim M=\dim N$。
\end{enumerate}
\end{proposition}

\begin{proof}
若 $F$ 开或闭,则由\cref{theorem:appA-embed-open-closed} 知它是拓扑嵌入,故它是光滑嵌入。(b) 或 (c) 都蕴涵 $F$ 闭:若 $F$ 固有,则由\cref{theorem:appA-proper-closed} 知它闭;若 $M$ 紧,则由闭映射引理知 $F$ 闭。最后,设 $M$ 边界为空且 $\dim M=\dim N$。则 $\mathrm{d}F_p$ 处处非奇异,而习题 4-2 表明 $F(M)\subseteq\operatorname{Int}N$。\cref{proposition:ch04-local-diffeo-equiv}(b) 表明 $F:M\to\operatorname{Int}N$ 是局部微分同胚,故它是开映射。于是 $F:M\to N$ 是开映射之复合 $M\to\operatorname{Int}N\hookrightarrow N$,故它是嵌入。
\end{proof}

\begin{example}\label{example:ch04-sphere-inclusion}
令 $\iota:S^n\hookrightarrow\mathbb{R}^{n+1}$ 为包含映射。我们在\cref{example:ch01-sphere-charts}(d) 中通过计算其关于图坐标的坐标表示证明了 $\iota$ 光滑。在同一坐标下容易验证其微分在每点单射,故它是单射光滑浸入。由于 $S^n$ 紧,由\cref{proposition:ch04-injective-immersion-embed} 知 $\iota$ 是光滑嵌入。
\end{example}

\begin{exercise}\label{exercise:ch04-embedding-not-open-closed}
举一个既不是开映射也不是闭映射的光滑嵌入的例子。
\end{exercise}

\begin{theorem}[Local Embedding Theorem]\label{theorem:ch04-local-embedding}
设 $M$ 与 $N$ 为带边或不带边的光滑流形, $F:M\to N$ 为光滑映射。则 $F$ 是光滑浸入,当且仅当 $M$ 中每一点都有邻域 $U\subseteq M$,使 $F|_U:U\to N$ 是光滑嵌入。
\end{theorem}

\begin{proof}
一个方向是显然的:若每一点都有邻域使 $F$ 在其上是光滑嵌入,则 $F$ 处处满秩,故它是光滑浸入。

反之,设 $F$ 是光滑浸入,并令 $p\in M$。我们先证明 $p$ 有一邻域使 $F$ 在其上是单射。若 $F(p)\notin\partial N$,则秩定理(当 $p\notin\partial M$)或\cref{theorem:ch04-immersion-boundary}(当 $p\in\partial M$)都蕴涵存在 $p$ 的邻域 $U_1$,使 $F$ 在其上具有形式\eqref{align:ch04-rank-immersion} 的坐标表示。由该式可知 $F|_{U_1}$ 单射。另一方面,设 $F(p)\in\partial N$,并任取以 $F(p)$ 为中心的 $N$ 的光滑边界坐标卡 $(W,\psi)$。令 $U_0=F^{-1}(W)$,它是 $p$ 的邻域,并令 $\kappa:\mathbb{H}^n\hookrightarrow\mathbb{R}^n$ 为包含映射,则可把前述论证应用于复合映射 $\kappa\circ\psi\circ F|_{U_0}:U_0\to\mathbb{R}^n$,以证明 $p$ 有邻域 $U_1\subseteq U_0$ 使 $\kappa\circ\psi\circ F|_{U_1}$ 单射,由此即得 $F|_{U_1}$ 单射。

现令 $p\in M$ 任意,并令 $U_1$ 为 $p$ 的邻域使 $F$ 在其上单射。存在 $p$ 的一个预紧邻域 $U$ 使 $\overline{U}\subseteq U_1$。$F$ 在 $\overline{U}$ 上的限制是定义域紧的单射连续映射,故由闭映射引理知它是拓扑嵌入。由于拓扑嵌入的任一限制仍是拓扑嵌入, $F|_U$ 既是拓扑嵌入又是光滑浸入,因而是光滑嵌入。
\end{proof}

\cref{theorem:ch04-local-embedding} 为一种对任意拓扑空间都有意义的浸入概念指明了方向:若 $X$ 与 $Y$ 是拓扑空间,称连续映射 $F:X\to Y$ 为\textbf{拓扑浸入},如果 $X$ 的每一点都有邻域 $U$ 使 $F|_U$ 是拓扑嵌入。于是,每个光滑浸入都是拓扑浸入;但正如嵌入的情形一样,恰好光滑的拓扑浸入未必是光滑浸入(参见\cref{example:ch04-smooth-topological-embedding})。

\end{document}